% ------------------------------- Περίληψη --------------------------------
\chapter*{Περίληψη}
\addcontentsline{toc}{chapter}{Περίληψη}
\thispagestyle{plain}

Οι περιφερειακές συσκευές \en{USB} αποτελούν έναν από τους πιο ευρέως
αξιοποιήσιμους φορείς επίθεσης με φυσική πρόσβαση. Μια συσκευή που δηλώνει ότι
είναι «στικάκι μνήμης» μπορεί να εμφανίζεται ταυτόχρονα στο λειτουργικό σύστημα ως
πληκτρολόγιο και να πληκτρολογεί εντολές μέσα σε λίγα χιλιοστά του δευτερολέπτου
(επιθέσεις τύπου \en{BadUSB}). Στο \en{Linux}, ο πυρήνας προσδένει οδηγό σε κάθε
νέα συσκευή αμέσως μόλις αυτή εμφανιστεί, χωρίς καμία ερώτηση προς τον χρήστη.

Η παρούσα εργασία παρουσιάζει τη σχεδίαση, υλοποίηση και αξιολόγηση του
\Probolos, μιας «πύλης αποδοχής» (\en{admission gate}) συσκευών \en{USB} με
προεπιλεγμένη άρνηση για το \en{Linux}. Το \Probolos{} αξιοποιεί τη διεπαφή
εξουσιοδότησης του πυρήνα (\code{authorized\_default}) ώστε κάθε νέα συσκευή να
εμφανίζεται \emph{απενεργοποιημένη}. Όσο η συσκευή αναμένει, εξετάζεται σε έως
τέσσερα στάδια: (α) ταυτότητα, από τους περιγραφείς της· (β) συνέπεια, με
21 κανόνες που εντοπίζουν αντιφατικούς ισχυρισμούς, όπως συσκευή αποθήκευσης που
δηλώνει και πληκτρολόγιο· (γ) συμπεριφορά, με σύντομη ενεργοποίηση της συσκευής
υπό αποκλειστική δέσμευση της εισόδου της (\code{EVIOCGRAB}) και ανάλυση του ρυθμού
πληκτρολόγησης· και (δ) αποθήκευση, με ανάγνωση μόνο-για-ανάγνωση του πίνακα
διαμερισμάτων χωρίς προσάρτηση. Η τελική απόφαση ανήκει πάντα στον άνθρωπο, μέσω
τερματικού ή διαλόγων επιφάνειας εργασίας με κλιμακούμενη «τριβή» ανάλογα με τον
κίνδυνο. Το σύστημα διατηρεί ιστορικό συσκευών με αποτυπώματα των περιγραφέων
για τον εντοπισμό «μετάλλαξης» (\en{descriptor drift}), εφαρμόζει διαχωρισμό
προνομίων ανάμεσα σε μια μικρή διεργασία \en{root} και έναν μη προνομιούχο
αναλυτή, και εγγυάται ότι η πύλη επανανοίγει σε κάθε έξοδο, ώστε ο χρήστης να μην
κλειδωθεί ποτέ έξω από το σύστημά του.

Η αξιολόγηση πραγματοποιήθηκε με εξομοιωμένες συσκευές \en{USB}
(\en{dummy\_hcd} και \en{raw\_gadget}), με σουίτα 1413 αυτοματοποιημένων ελέγχων,
ελέγχους βασισμένους σε ιδιότητες, ασαφή έλεγχο (\en{fuzzing}) καθοδηγούμενο
από κάλυψη και επανειλημμένους ελέγχους κώδικα (\en{audits}). Καταγράφηκαν και
διορθώθηκαν 106 ελαττώματα, εκ των οποίων 4 κρίσιμης προτεραιότητας (P0). Η
εργασία καταλήγει ότι μια πύλη προεπιλεγμένης άρνησης, συνδυασμένη με ελέγχους
συνέπειας, μειώνει ουσιαστικά την επιφάνεια επίθεσης \en{BadUSB} με χαμηλό
κόστος χρηστικότητας, και καταγράφει με ειλικρίνεια τα όρια της προσέγγισης.

\vspace{0.8cm}
\noindent\textbf{Λέξεις-κλειδιά:} ασφάλεια \en{USB}, \en{BadUSB}, έγχυση
πληκτρολογήσεων, \en{Linux}, \en{sysfs}, εξουσιοδότηση συσκευών, διαχωρισμός
προνομίων, ασφάλεια περιφερειακών, ανίχνευση ανωμαλιών.
